\chapter{Doctoral Entrance Problems}
\label{ch:problems}

The problems below are reproduced in French and routed by individual mathematical content.
\newpage

\section{Université des Sciences et de la Technologie Houari Boumediène (USTHB)}
\begin{problem}{(USTHB) — Épreuve N°01}{pde:usthb-2026-e1-p1}
Soit $\Omega\subset\mathbb R^2$ un ouvert borné de frontière régulière. On considère le problème variationnel
$$\int_\Omega\nabla u\cdot\nabla v\,dxdy+\int_\Omega(\nabla u\cdot F)v\,dxdy=\int_\Omega f\cdot\nabla v\,dxdy,\qquad\forall v\in H_0^1(\Omega),$$
avec $f\in(L^2(\Omega))^2$, $F\in(L^\infty(\Omega))^2$ et $\operatorname{div}(F)=0$. Montrer que $\nabla u\cdot F=\operatorname{div}(uF)$, que $\int_\Omega u(F\cdot\nabla u)=0$, établir l'existence et l'unicité de $u\in H_0^1(\Omega)$, puis déterminer le problème aux limites satisfait par $u$ lorsque $u\in H^2(\Omega)$ et $\operatorname{div}(f)\in L^2(\Omega)$.
\source{\emph{\small Université des Sciences et de la Technologie Houari Boumediène (USTHB), 2026, Épreuve~n\textsuperscript{o}~01.}}
\end{problem}
\newpage

\section{Université Mohamed Boudiaf --- M'Sila}
\begin{problem}{(M'Sila) — Épreuve N°01}{pde:msila-2025-e1-p1}
Soit $\Omega$ un ouvert borné de $\mathbb R^2$ de frontière régulière. On considère
$$\int_\Omega\nabla u\cdot\nabla v+\int_\Omega u(F\cdot\nabla v)=\int_\Omega f\cdot\nabla v,\qquad\forall v\in H_0^1(\Omega),$$
avec $f\in(L^2(\Omega))^2$, $F\in(L^\infty(\Omega))^2$ et $\operatorname{div}F=0$. Montrer l'existence et l'unicité de $u\in H_0^1(\Omega)$ et, sous $u\in H^2(\Omega)$ et $\operatorname{div}f\in L^2(\Omega)$, déterminer le problème aux limites associé.
\source{\emph{\small Université Mohamed Boudiaf - M'Sila, 2025, Épreuve~n\textsuperscript{o}~01.}}
\end{problem}
\newpage

\section{Université Chadli Bendjedid --- El Tarf}
\begin{problem}{(El Tarf) — Épreuve N°01}{pde:eltarf-2025-e1-p1}
On considère l'équation aux dérivées partielles du second ordre
$$y^2u_{xx}-e^{2x}u_{yy}=1,\qquad(x,y)\in\mathbb R\times\mathbb R^*.$$
Déterminer la nature de cette équation et ses courbes caractéristiques.
\source{\emph{\small Université Chadli Bendjedid - El Tarf, 2025, Épreuve~n\textsuperscript{o}~01.}}
\end{problem}
\newpage

\begin{problem}{(El Tarf) — Épreuve N°01}{pde:eltarf-2025-e1-p2}
On considère le problème de la chaleur sur $]0,\pi[$ :
$$\begin{cases}u_t-u_{xx}=0,\\u(0,t)=u(\pi,t)=0,\\u(x,0)=2\cos\!\bigl(3(x+\tfrac\pi6)\bigr).\end{cases}$$
Simplifier la donnée initiale et résoudre le problème par séparation des variables.
\source{\emph{\small Université Chadli Bendjedid - El Tarf, 2025, Épreuve~n\textsuperscript{o}~01.}}
\end{problem}
